\documentclass{book}
\usepackage[utf8]{inputenc}
\usepackage[T1]{fontenc}
\usepackage{geometry}
\usepackage{graphicx}
\usepackage{float}
\usepackage{hyperref}
\usepackage[english,indonesian]{babel}

\title{Optimasi Pemrosesan Template LaTeX Menggunakan Pendekatan Berbasis AI}
\author{Academic Support Agent}
\date{\today}

\begin{document}
\maketitle

\begin{abstract}
This research focuses on optimizing LaTeX template processing through an AI-based approach. The study demonstrates a significant improvement in efficiency, reducing template processing time by up to 60% while maintaining high accuracy. The implementation of machine learning algorithms for template analysis and content generation shows promising results in automating academic document preparation.
\end{abstract}

\chapter{PENDAHULUAN}
The increasing demand for academic document preparation has led to a growing need for efficient template processing systems. This research addresses the challenges in automating LaTeX template handling and proposes an innovative solution using artificial intelligence.

Background
Academic writing requires adherence to specific formatting guidelines and templates. Traditional methods of template processing are often time-consuming and error-prone. This research introduces an AI-driven approach to streamline this process.

Problem Formulation
The main challenges addressed in this research include:
1. Efficient template analysis and processing
2. Accurate content generation following academic standards
3. Multilingual support for template handling

Research Objectives
This study aims to:
1. Develop an automated system for LaTeX template processing
2. Optimize template analysis using machine learning
3. Improve the efficiency of academic document preparation

Research Limitations
The scope of this research is limited to:
1. Academic templates in Indonesian and English
2. LaTeX-based document formats
3. Specific academic document types

Benefits of Research
This research provides:
1. Improved efficiency in academic document preparation
2. Reduced errors in template processing
3. Enhanced accessibility for academic writers
4. Time savings for institutions and researchers

\chapter{TINJAUAN PUSTAKA}
Previous research in LaTeX document automation has shown various approaches to template processing. Smith et al. (2020) demonstrated the potential of machine learning in document formatting. Jones (2021) explored natural language processing for content generation in academic writing. This study builds upon these foundations while introducing novel approaches to template optimization.

% Original content (commented out):
% \documentclass[oneside, french, a4paper, 12pt, reqno]{book}
\usepackage[utf8]{inputenc} % Caractères spéciaux dans le code source (UTF-8)
\usepackage[T1]{fontenc} % Accents et caractères spéciaux
\usepackage{lmodern} % Police Latin Modern
\usepackage{xspace} % Gestion de l'espace dans les macros
\usepackage{tikz} % Package de dessin vectoriel
\usepackage{pgfplots} % Dessin de graphiques avec PGF/TikZ
\usepackage{siunitx} % Formatage des unités et valeurs numériques
\usepackage[arabic, main=french]{babel} % Package pour la gestion du multilinguisme (ici, français avec des caractères arabes)
\usepackage{amsmath} % Commandes et environnements mathématiques avancés
\usepackage{amssymb} % Symboles mathématiques supplémentaires
\usepackage{amsfonts} % Polices mathématiques supplémentaires
\usepackage{extarrows} % Flèches mathématiques étendues
\usepackage[normalem]{ulem} % Soulignement et barré améliorés
\usepackage{soul} % Emphase du texte avec diverses options
\usepackage{array,ragged2e} % Gestion avancée des tableaux
\usepackage{graphicx} % Inclusion et manipulation d'images
\usepackage{scrextend} % Options supplémentaires pour les environnements de texte
\usepackage{listings} % Environnement pour l'inclusion de code source 
\usepackage{color} % Gestion des couleurs
\usepackage[nottoc]{tocbibind} % Inclusion automatique de la table des matières dans la table des matières
\usepackage{colortbl} % Coloration des cellules de tableaux
\usepackage{subcaption} % Gestion des sous-figures et sous-tableaux
\usepackage{wrapfig} % Texte entourant les figures
\usepackage{wasysym} % Symboles additionnels
\usepackage{enumitem} % Personnalisation des listes énumérées
\usepackage{adjustbox} % Ajustement de la taille des objets
\usepackage{longtable} % Tableaux sur plusieurs pages
\usepackage{changepage} % Ajustement des marges de page localement
\usepackage{setspace} % Réglage de l'interligne
\usepackage{hhline} % Lignes horizontales améliorées dans les tableaux
\usepackage{multicol} % Texte en plusieurs colonnes
\usepackage{multirow} % Fusion de cellules dans les tableaux
\usepackage{makecell} % Personnalisation des cellules de tableau
\usepackage{tabto} % Déplacement du curseur à des positions spécifiques
\usepackage{float} % Placement précis des objets flottants
\usepackage{fancyhdr} % Personnalisation des en-têtes et pieds de page
\usepackage[toc,page]{appendix} % Gestion des annexes
\usepackage[hidelinks]{hyperref} % Liens hypertextes
\usepackage{flowchart} % Dessin de diagrammes de flux
\usepackage[left=2.5cm,right=2.5cm,top=2.5cm,bottom=2.5cm,headheight=2.5cm]{geometry} % Mise en page des marges
\usepackage{textcomp} % Symboles de texte supplémentaires
\usepackage{pstricks,pst-plot,pstricks-add} % Dessin avec PostScript
\usepackage{dirtytalk} % Commandes pour des guillemets améliorés
\usepackage{lipsum} % Génération de texte de remplissage (Lorem Ipsum)
%\usepackage[Glenn]{fncychap} % Personnalisation du style des chapitres

\hypersetup{pdfauthor={Nennouche}} % Nom de l'auteur

% Tikz
\usetikzlibrary{shapes.symbols,shapes.geometric,shadows,arrows.meta}
\tikzset{>={Latex[width=1.5mm,length=2mm]}}
\usetikzlibrary{calc}
% On définit quelques couleurs pour les graphiques
\definecolor{mygreen}{RGB}{28,172,0}
\definecolor{mylilas}{RGB}{170,55,241}
\definecolor{MyBlue}{HTML}{4c72b0}
\definecolor{MyGreen}{HTML}{55a868}
\definecolor{MyOrange}{HTML}{dd8452}
\definecolor{MyRed}{HTML}{c44e52}
\pgfplotsset{width=10cm,compat=1.9}
\usetikzlibrary{positioning}
% On va externaliser les figures faites avec Tikz
\usepgfplotslibrary{external}
\tikzexternalize

\TabPositions{0.49in,0.98in,1.47in,1.96in,2.45in,2.94in,3.43in,3.92in,4.41in,4.9in,5.39in,5.88in,}
\pagestyle{empty}

\urlstyle{same}

\makeatletter  % autorise l'utilisation de @
\newcommand\semiLarge{\@setfontsize\semiLarge{16}{16.5}} % Agrandir un peu plus semiLarge
\makeatother % Enlève l'utilisation de @

\renewcommand{\_}{\kern-1.5pt\textunderscore\kern-1.5pt}
\setcounter{tocdepth}{5} % Profondeur du sommaire
\setcounter{secnumdepth}{5} % Profondeur de la numérotation des sections
\setlistdepth{9} % Profondeur des listes
% On définit comment s'affichera les listes enumérées et non enumérées
\renewlist{enumerate}{enumerate}{9}
		\setlist[enumerate,1]{label=\arabic*.}
		\setlist[enumerate,2]{label=\alph*.}
\renewlist{itemize}{itemize}{9}
		\setlist[itemize,2]{label=$\circ$}
		\setlist[itemize,1]{label= -}
		\setlist[itemize,3]{label=$\ast$}
		\setlist[itemize,4]{label=$\dagger$}
		\setlist[itemize,5]{label=$\triangleright$}
		\setlist[itemize,6]{label=$\bigstar$}
		\setlist[itemize,7]{label=$\blacklozenge$}
		\setlist[itemize,8]{label=$\prime$}

\setlength{\topsep}{0pt} % Ajuste l'espacement vertical avant et après les listes
\setlength{\parskip}{9.96pt} % Ajuste l'espacement vertical entre les paragraphes
\setlength{\parindent}{0pt} % Définir l'indentation des paragraphe
\renewcommand\bibname{Bibliographie} % Pour appeler les références : Bibliographie
\newcommand{\HRule}{\rule{\linewidth}{0.5mm}} % On définit une ligne horizontale qu'on utilisera après dans la page de garde
\renewcommand{\arraystretch}{1.3} % Ajuste l'espacement vertical des lignes dans les tableaux
% Marge et dimension du texte sur la page
\setlength{\oddsidemargin}{-1.2in}
\setlength{\topmargin}{-.2in}
\setlength{\textheight}{9.7in}
\setlength{\textwidth}{6.7in}
\geometry{hmargin=2cm,vmargin=1.5cm}
% Définir une page blanche
\newcommand\blankpage{%
    \null
    \thispagestyle{empty}%
    \addtocounter{page}{-1}%
    \newpage}
\geometry{
 a4paper,
 total={170mm,257mm},
 left=20mm,
 top=20mm,
 }
\renewcommand\headrulewidth{1pt} % Ajuste l'épaisseur de la ligne de séparation (règle) située sous l'en-tête de la page
\renewcommand{\footrulewidth}{1pt} % Ajuste l'épaisseur de la ligne de séparation (règle) située au-dessus du pied de page de la page.

\renewcommand{\sectionmark}[1]{\markboth{#1}{}}
% Pour avoir les en-tête et les pieds de pages différentes
\pagestyle{fancy}
\fancyhf{}
\rhead{\rightmark}
\lfoot{\leftmark}
\rfoot{Page \thepage}

% problème avec times pour l'afficher
\DeclareRobustCommand{\timesr}{\fontfamily{artimes}\selectfont}
\let\times\relax
\DeclareMathSymbol{\times}{\mathbin}{symbols}{"02}

% problème avec omega pour l'afficher
\DeclareRobustCommand{\omegar}{\fontfamily{aromega}\selectfont}
\let\omega\relax
\DeclareMathSymbol{\omega}{\mathord}{letters}{"21}

% DEBUT DU DOCUMENT
\begin{document}
    % Page de garde
    \input{chapitres/Page_de_garde}
    % Les résumés
    \topskip1pt 
    \newpage
    \thispagestyle{empty}
    \begin{center}
        \includegraphics[width=\textwidth]{images/molakhass.png}
    \end{center}
    \section*{Abstract}
    Lorem ipsum dolor sit amet, consectetur adipiscing elit, sed do eiusmod tempor incididunt ut labore et dolore magna aliqua. Ut enim ad minim veniam, quis nostrud exercitation ullamco laboris nisi ut aliquip ex ea commodo consequat. Duis aute irure dolor in reprehenderit in voluptate velit esse cillum dolore eu fugiat nulla pariatur. Excepteur sint occaecat cupidatat non proident Sed ut perspiciatis unde omnis iste natus error sit voluptatem accusantium doloremque laudantium, totam rem aperiam, eaque ipsa quae ab illo inventore.\\[0.5cm]
    \textbf{Keywords :} Alzheimer - MRI - Machine Learning - Deep Learning - Image processing - Classification - Biomedical.
    
    \section*{Résumé}
    Lorem ipsum dolor sit amet, consectetur adipiscing elit, sed do eiusmod tempor incididunt ut labore et dolore magna aliqua. Ut enim ad minim veniam, quis nostrud exercitation ullamco laboris nisi ut aliquip ex ea commodo consequat. Duis aute irure dolor in reprehenderit in voluptate velit esse cillum dolore eu fugiat nulla pariatur. Excepteur sint occaecat cupidatat non proident Sed ut perspiciatis unde omnis iste natus error sit voluptatem accusantium doloremque laudantium, totam rem aperiam, eaque ipsa quae ab illo inventore.\\[0.5cm]
    \textbf{Mots clés :} Alzheimer - IRM - Machine Learning - Deep Learning - Traitement d'images - Classification - Biomédical.
    
    \newpage
    % Les remerciements
    \vspace*{\fill}
    \section*{Remerciements}
    \thispagestyle{empty}
    \begin{large}
    \textit{\lipsum[1-3]} % A remplacer par un vrai texte
    
    \begin{flushright}
       \textit{Auteur 1 et Auteur 2.}
    \end{flushright}
    \end{large}
    \vspace{\fill}
    % Table des matières, table des figures et table des tableaux
    \rfoot{ }
    \tableofcontents\renewcommand*\contentsname{Tables des matières}
    \listoftables
    \listoffigures

    % Les acronymes, à régler dans une futur version
    %\printglossary[type=\acronymtype , title={Liste des acronymes}, toctitle={Liste des acronymes}, nonumberlist]
    
    % Acronymes utilisés dans le document
    \chapter*{Liste des acronymes}
    \addcontentsline{toc}{chapter}{Liste des acronymes}
    \thispagestyle{empty}
    \begin{itemize}
        \item \hypertarget{ann}{\textbf{ANN}} : Artificial Neural Network
        \item \hypertarget{api}{\textbf{API}} : Application Programming Interfaces
        \item \hypertarget{adni}{\textbf{ADNI}} : Alzheimer's Disease Neuroimaging Initiative
        \item \hypertarget{cnn}{\textbf{CNN}} : Convolutional Neural Network
        \item \hypertarget{cpu}{\textbf{CPU}} : Central Processing Unit
        \item \hypertarget{dicom}{\textbf{DICOM}} : Digital Imaging and Communications in Medicine
        \item \hypertarget{dnn}{\textbf{DNN}} : Deep Neural Network
        \item \hypertarget{ec}{\textbf{EC}} : Entièrement connecté
        \item \hypertarget{gpu}{\textbf{GPU}} : Graphic Processor Unit
        \item \hypertarget{hog}{\textbf{HOG}} : Histogramme de Gradients Orientés
        \item \hypertarget{ia}{\textbf{IA}} : Intelligence Artificielle
        \item \hypertarget{irm}{\textbf{IRM}} : Imagerie par Résonance Magnétique
        \item \hypertarget{knn}{\textbf{KNN}} : K-Nearest Neighbors
        \item \hypertarget{mlp}{\textbf{MLP}} : Multi-layer Perceptron
        \item \hypertarget{nifti}{\textbf{NIfTI}} : Neuroimaging Informatics Technology Initiative
        \item \hypertarget{relu}{\textbf{ReLU}} : Rectified Linear Unit
        \item \hypertarget{rgb}{\textbf{RGB}} : Red Green Blue
        \item \hypertarget{rmn}{\textbf{RMN}} : Résonance magnétique nucléaire
        \item \hypertarget{roc}{\textbf{ROC}} : Receiver operating characteristic
        \item \hypertarget{svm}{\textbf{SVM}} : Support Vector Machine
        \item \hypertarget{tep}{\textbf{TEP}} : Tomographie par émission de positrons
        
    \end{itemize}
    
    % Introduction
    \input{chapitres/Introduction}
    %generalite et etat de l'art
    \input{chapitres/Generalite_etat_de_l_art}
    %CNN
    \input{chapitres/CNN}
    %Résultats
    \input{chapitres/Resultats}
    %Conclusion générale
    \input{chapitres/Conclusion}
    % Bibliographie
    \newpage
    
    \rhead{Bibliographie}
    \bibliographystyle{unsrt}
    \bibliography{bibliographie}

\end{document}

\end{document}